\documentclass[12pt]{article}
\usepackage[T1]{fontenc}
\usepackage[utf8]{inputenc}
\usepackage{lmodern}
\usepackage{textcomp}
\usepackage{enumitem}
\usepackage{geometry}
\usepackage{graphicx}
\usepackage{amsmath, amssymb}
\geometry{margin=1in}
\begin{document}
\begin{center}
Sociology - Graduate Level Assessment 11 \
Total Marks: 40 \
Answer all questions.
\end{center}

\section*{Multiple Choice (10 marks)}
\begin{enumerate}[label=\arabic*.]
\item Bernstein thought that using 'restricted codes' of language disadvantaged pupils at school because:
\begin{enumerate}[label=(\alph*)]
    \item this pattern of speech made them the target of bullying
    \item they referred to explicit, context independent meanings
    \item they prevented children from communicating outside of their peer groups
    \item they involved short, simple sentences with a small vocabulary
\end{enumerate}
\item Mulholland (1998) argued that privatization changed the relationship between companies and managers in that:
\begin{enumerate}[label=(\alph*)]
    \item flexibility was reduced by the introduction of detailed daily worksheets
    \item the state had greater control than managers over production processes
    \item ownership was transferred from small shareholders to senior managers
    \item employment depended on performance rather than trust and commitment
\end{enumerate}
\item How is terrorism different from the types of crime described by the Chicago School?
\begin{enumerate}[label=(\alph*)]
    \item it is committed on a larger, often global, scale, and is well organized
    \item it is associated with political conflict between states and their citizens
    \item it can have far-reaching effects upon international relations
    \item all of the above
\end{enumerate}
\item The middle classes that developed over the nineteenth century were:
\begin{enumerate}[label=(\alph*)]
    \item an urban set, involved in civic bodies and voluntary associations
    \item too diverse to have a strong sense of class consciousness
    \item often involved in 'white collar' work
    \item all of the above
\end{enumerate}
\item The capitalist class of the mid-twentieth century were said to join the upper class because they:
\begin{enumerate}[label=(\alph*)]
    \item participated in the same leisure pursuits and events of the 'social calendar'
    \item emulated the lifestyle and cultural values of the traditional aristocracy
    \item owned companies and financial assets that generated wealth through corporations
    \item had direct, personal ownerships of land and businesses as physical assets
\end{enumerate}
\end{enumerate}

\section*{True / False (10 marks)}
\textit{Answer True or False}
\begin{enumerate}[label=\arabic*.]
\setcounter{enumi}{5}
\item True or False: In stage 3 of the 'health transition', acute infectious diseases such as typhus and cholera remain the main causes of illness and death.
\item True or False: The rates of violent crime are similar for men and women, suggesting equal levels of aggression among genders.
\item True or False: The social democratic model emphasizes individualistic self-reliance and supports deregulation of wages and prices.
\item True or False: The growing fear of HIV and AIDS, influenced by the New Right, did not arise until after the sexual revolution of the 1960 s.
\item True or False: A document is deemed 'authentic' if it represents all similar documents that did not survive.
\end{enumerate}

\section*{Long Form Response (20 marks)}
\textit{Answer each question with a clear and complete explanation.}
\begin{enumerate}[label=\arabic*.]
\setcounter{enumi}{10}
\item Discuss the implications of globalization on the power dynamics of nation-states, specifically addressing which aspects of globalization challenge traditional state authority.
\item Discuss Wallerstein's concept of a 'world system' in the context of the capitalist world economy, emphasizing the significance of a division of labor across diverse ethnic and cultural groups.
\end{enumerate}

\vfill
\noindent\textit{End of Paper}
\end{document}